% Generated draft: requires mathematical and source review.
\documentclass[11pt]{article}
\usepackage[margin=25mm]{geometry}
\usepackage{fontspec}
\setmainfont{Latin Modern Roman}
\usepackage{amsmath,amssymb,mathtools,unicode-math}
\setmathfont{Latin Modern Math}
\usepackage{enumitem,needspace,xurl,hyperref}
\setlength{\parindent}{0pt}
\setlength{\parskip}{5pt}
\setlength{\emergencystretch}{3em}
\newcommand{\N}{\mathbb{N}}
\newcommand{\Z}{\mathbb{Z}}
\newcommand{\Q}{\mathbb{Q}}
\newcommand{\R}{\mathbb{R}}
\newcommand{\C}{\mathbb{C}}
\newcommand{\F}{\mathbb{F}}
\newcommand{\A}{\mathbb{A}}
\newcommand{\PP}{\mathbb{P}}
\newcommand{\E}{\mathbb{E}}
\newcommand{\eps}{\varepsilon}
\newcommand{\dd}{\,\mathrm{d}}
\begin{document}
\section*{UnsolvedMath ID: 800686}
\textbf{Source problem code:} catalogue entry
\subsection{Definitions and mathematical statement}
Let \(M\) be a complex manifold of complex dimension \(n\), homogeneous under an action of a Lie group \(G\); thus the action on \(M\) is transitive. Write \(\PP^n=\PP^n(\C)\) and set
\[
\mathscr M=\operatorname{Hol}(M,\PP^n).
\]
Equip \(\mathscr M\) with the topology of uniform convergence on compact subsets of \(M\), using, for example, the Fubini--Study distance \(d_{\mathrm{FS}}\) on \(\PP^n\).

The source states that \(G\) acts on \(\mathscr M\), without giving the action formula. Under the usual interpretation that each element of \(G\) acts on \(M\) by a biholomorphism, the induced action is
\[
(g\cdot f)(x)=f(g^{-1}\cdot x),
\qquad g\in G,\quad f\in\mathscr M,\quad x\in M.
\]
This interpretation requires that the transformations of \(M\) preserve its complex structure, a hypothesis not explicitly stated in the supplied record.

\begin{enumerate}
\item[(i)] Determine conditions on \(M\), \(G\), and the homogeneous action under which the induced action on \(\mathscr M\) admits a dense orbit. For the precomposition action above, the property to be characterized is
\[
\exists f\in\mathscr M\quad
\overline{\{g\cdot f:g\in G\}}=\mathscr M.
\]
Equivalently, there exists \(f\in\mathscr M\) such that for every \(h\in\mathscr M\), every nonempty compact subset \(K\subseteq M\), and every \(\varepsilon>0\), there is \(g\in G\) satisfying
\[
\sup_{x\in K}d_{\mathrm{FS}}\bigl(f(g^{-1}\cdot x),h(x)\bigr)<\varepsilon.
\]
\item[(ii)] Extend the construction attributed in the source to B. Weiss to this setting. The supplied record does not describe that construction or give its reference, hypotheses, or intended conclusion. Consequently, this part remains an unresolved request for an extension rather than a fully specified mathematical assertion.
\end{enumerate}
\subsection{Short English statement}
For a homogeneous complex manifold \(M\) of complex dimension \(n\), determine when the induced action of \(G\) on holomorphic maps \(M\to\PP^n\) has a dense orbit in the topology of uniform convergence on compact sets. Also extend the construction attributed to B. Weiss, whose details are absent from the supplied record. The displayed action uses the stated interpretation that \(G\) acts by biholomorphisms.
\subsection{Sources}
\begin{itemize}
\item Local UnsolvedMath registry, supplied complete source record: permanent ID \(800686\), exact source problem code \texttt{catalogue entry}, AMR Open Problem Lists, title ``Holomorphic Foliation Theory --- Problem 13.'' The registry attributes the problem to Dinh and Sibony, with proposed year 2018. Its \texttt{open} status is a supplied registry label, not an independently verified mathematical or literature assessment.
\item Dinh and Sibony, \emph{Some Open Problems on Holomorphic Foliation Theory} (2018), Problem 13. Supplied source URL: \url{https://arxiv.org/abs/1807.05640}. The supplied provenance records access on 2026-10-16 and extraction by \texttt{arxiv\_latex}; these are supplied provenance details, not a claim of independent source access or verification for this draft.
\item B. Weiss is credited in part (ii) with a construction to be extended. No title, bibliographic reference, or description of that construction is supplied.
\end{itemize}
\end{document}
